\documentclass[11pt]{article}
\ifdefined\pdfmapfile\pdfmapfile{+lm.map}\pdfmapfile{+cm.map}\pdfmapfile{+symbols.map}\fi
\usepackage[T1]{fontenc}
\usepackage{lmodern}
\usepackage[a4paper,margin=25mm]{geometry}
\usepackage{amsmath,amssymb,amsthm,booktabs,hyperref,listings,xcolor}
\hypersetup{colorlinks=true,urlcolor=blue,citecolor=blue,linkcolor=blue}
\newtheorem{theorem}{Theorem}
\newtheorem{lemma}[theorem]{Lemma}
\newtheorem{proposition}[theorem]{Proposition}
\newcommand{\CM}{\mathcal{CM}}
\newcommand{\R}{\mathbb R}
\newcommand{\fall}[2]{#1^{\underline{#2}}}
\lstset{basicstyle=\ttfamily\scriptsize,columns=fullflexible,keepspaces=true,
  breaklines=true,showstringspaces=false,frame=single,language=Python}
\title{Real powers of completely monotone functions\protect\\have sharp derivative threshold twelve}
\author{Research note}
\date{14 November 2026}
\begin{document}
\maketitle
\begin{abstract}
For every completely monotone function $f$ on $(0,\infty)$ and every real
$\alpha\ge1$, we prove that $f^\alpha$ satisfies the alternating derivative
inequalities through order twelve. The bound is sharp, as shown by an exact
order-thirteen counterexample. The proof reduces derivatives to Stieltjes
moment polynomials and supplies a finite certificate of nonnegativity using
rational Bernstein coefficients. All 2048 composition monomials are covered
by 64 groups and 145 terminal regions, with maximum subdivision depth twelve.
The complete exact verifier is included. This determines the boundary between
the order-seven preservation and order-thirteen obstruction documented by
van Haeringen in 1997. The argument is computer-assisted and is not a
proof-assistant formalization; external review is pending.
\end{abstract}

\section{Statement and historical scope}

A function $f\in C^\infty(0,\infty)$ is completely monotone if
$(-1)^jf^{(j)}(t)\ge0$ for every integer $j\ge0$ and every $t>0$.
Write $\CM^N$ for the analogous class with only $0\le j\le N$ required.

\begin{theorem}\label{thm:main}
If $f$ is completely monotone on $(0,\infty)$ and $\alpha\ge1$, then
\[
(-1)^j\frac{d^j}{dt^j}f(t)^\alpha\ge0
\qquad(t>0,\quad 0\le j\le12).
\]
\end{theorem}

Van Haeringen's 1997 paper \cite{vH} states in its published abstract that the
universal implication $f\in\CM\Rightarrow f^\alpha\in\CM^N$ for every
$\alpha>1$ holds for $N\le7$ and fails for $N\ge13$.
The order-seven inclusion is also explicitly quoted in Lemma~5.3(2) of the
2026 preprint \cite{SW}. These are primary sources for the older question and
its continued use. They are not an exhaustive priority search, and we do not
attribute to van Haeringen an additional named conjecture not verified in the
original source. The present result closes orders eight through twelve in
that stated range and proves that twelve is the sharp universal threshold.
It concerns powers of functions completely monotone to every order; it does
not assert the corresponding preservation theorem for an input known only
to lie in $\CM^{12}$.

\section{Reduction to a polynomial in Stieltjes moments}

If $f\equiv0$, the theorem is immediate. Otherwise Bernstein's representation
theorem \cite[Theorem~1.4]{Bernstein} gives a positive measure $\mu$ on $[0,\infty)$ such that
\[
f(t)=\int_0^\infty e^{-tx}\,d\mu(x)>0.
\]
The measure $\mu$ need not have finite mass. Fix $t_0>0$ and form the probability
measure
\[
d\nu(x)=\frac{e^{-t_0x}}{f(t_0)}\,d\mu(x),\qquad
m_j=\int x^j\,d\nu(x)=\frac{(-1)^jf^{(j)}(t_0)}{f(t_0)}.
\]
All its moments are finite: for each $j$ there is a finite constant $C$ such
that $x^je^{-t_0x}\le C e^{-t_0x/2}$, whose integral is $Cf(t_0/2)$.

Let $B_{n,k}$ denote the exponential partial Bell polynomials, specified by
\begin{equation}\label{eq:Bell}
B_{0,0}=1,\qquad
B_{n,k}=\sum_{i=1}^{n-k+1}\binom{n-1}{i-1}m_iB_{n-i,k-1},
\end{equation}
with the other boundary values $B_{n,0}=B_{0,k}=0$ for positive indices.
The finite chain rule yields
\begin{equation}\label{eq:derivative}
\frac{(-1)^n(f^\alpha)^{(n)}(t_0)}{f(t_0)^\alpha}
=D_n(\alpha;m):=\sum_{k=1}^n\fall{\alpha}{k}B_{n,k}(m_1,m_2,\ldots),
\end{equation}
where $\fall{\alpha}{k}=\alpha(\alpha-1)\cdots(\alpha-k+1)$.
Every monomial in $B_{n,k}$ has weighted degree $n$ in the moments, which
accounts for the common sign $(-1)^n$. Formula~\eqref{eq:derivative} holds for
noninteger $\alpha$ because the scalar function $x^\alpha$ is smooth at the
positive value $f(t_0)$. No infinite Taylor-series argument is needed.

The following classical parametrization is the point at which all moment
inequalities are incorporated, rather than merely a log-convexity condition.

\begin{lemma}[Stieltjes parametrization]\label{lem:Stieltjes}
For the moments of any probability measure on $[0,\infty)$ having all moments
finite, there are $q_1,q_2,\ldots\ge0$ such that, as formal power series,
\[
\sum_{j\ge0}m_jz^j=
\cfrac{1}{1-\cfrac{q_1z}{1-\cfrac{q_2z}{1-\cdots}}}.
\]
In particular, for $n\ge1$,
\begin{equation}\label{eq:compositions}
m_n=\sum_{\substack{e_1+\cdots+e_h=n\\e_i\ge1}}
\left[\prod_{i=1}^{h-1}\binom{e_i+e_{i+1}-1}{e_{i+1}}\right]
q_1^{e_1}\cdots q_h^{e_h}.
\end{equation}
\end{lemma}

The moment/continued-fraction equivalence is Theorem~1.1(c)--(d) of
P\'etr\'eolle--Sokal--Zhu \cite{PSZ}, which attributes it to Stieltjes (1894).
It includes nonnegative parameters and degenerate measures. It is explicitly
a formal-power-series statement; convergence of $\sum m_jz^j$ is unnecessary.
The weighted Dyck-path interpretation is due to Flajolet \cite{Flajolet}.
For completeness, a Dyck path corresponds to an ordered rooted tree. If $e_i$
is the number of vertices at level $i$, the $e_{i+1}$ ordered children can be
distributed among $e_i$ ordered parents in
$\binom{e_i+e_{i+1}-1}{e_{i+1}}$ ways. Multiplying these independent counts
proves~\eqref{eq:compositions}.

There is also a direct continuity safeguard for degeneracies. Replace $\nu$ by
$(1-\varepsilon)\nu+\varepsilon\operatorname{Exp}(1)$. All ordinary and shifted
Hankel forms of this measure are strictly positive, because the exponential
component integrates every nonzero polynomial square, with or without an
extra factor $x$, to a positive number. Its moments converge to those of $\nu$.
Any polynomial inequality proved for the strict case therefore passes to
$\nu$. This covers point masses, an atom at zero, and vanishing Hankel
determinants without division by a determinant equal to zero.

\section{The exact order-twelve certificate}

Put $\alpha=1+\beta$ and initially restrict $0\le\beta\le1$. Substitute
\eqref{eq:compositions} into
\[
P_n(\beta;q)=\frac{D_n(1+\beta;m(q))}{1+\beta}.
\]
This is a polynomial with integer coefficients, homogeneous of degree $n$ in
$q$. Its support consists of the $2^{n-1}$ compositions of $n$, padded by
zeros. The exact verifier checks this support count at $n=12$. Equations
\eqref{eq:Bell}--\eqref{eq:compositions} specify every coefficient by finite
integer arithmetic. An equivalent first-return recurrence, used by the
verifier, is
\begin{equation}\label{eq:firstreturn}
M_{h,0}=1,\qquad
M_{h,n}=q_{h+1}\sum_{j=0}^{n-1}M_{h,j}M_{h+1,n-1-j},\qquad m_n=M_{0,n}.
\end{equation}
Here $M_{h,n}$ counts weighted Dyck paths based at height $h$. Both
constructions have also been independently compared coefficient by coefficient
for all 2048 monomials of $D_{12}$.

Group monomials of $P_{12}$ by their tail $\tau=(e_7,\ldots,e_{12})$.
There are 64 groups. We give separate certificates for the 63 nonzero tails
and the one zero tail.

\subsection{The 63 nonzero-tail groups: cubes}

For each fixed nonzero tail, factor out its monomial
$q_7^{e_7}\cdots q_{12}^{e_{12}}$ and the componentwise minimum powers
$q_1^{a_1}\cdots q_6^{a_6}$ appearing in that group. The remaining polynomial
has the form
\[
H_\tau(\beta;q_1,\ldots,q_6)=
\sum_{|u|=d}c_u(\beta)q_1^{u_1}\cdots q_6^{u_6}.
\]
All removed factors are nonnegative. Homogeneity reduces nonnegativity to
$q_1+\cdots+q_6=1$; the zero-sum case follows by continuity (or is constant
when $d=0$). Parametrize this simplex surjectively by
\begin{equation}\label{eq:stick}
q_j=x_j\prod_{h<j}(1-x_h)\quad(1\le j\le5),\qquad
q_6=\prod_{h=1}^5(1-x_h),\qquad 0\le x_j\le1.
\end{equation}
This covers all boundary faces as well as the interior.

Write $c_u(\beta)=\sum_\ell c_{u,\ell}\beta^\ell$ and
$K=\max_u\deg c_u$. Use the tensor Bernstein basis of degrees
$(K,d,d,d,d,d)$ in $(\beta,x_1,\ldots,x_5)$. Put
$s_j=\sum_{h<j}u_h$. For $0\le i\le K$ and $0\le v_j\le s_j$, the monomial
$c_u(\beta)q^u$ contributes at index $(i,u_1+v_1,\ldots,u_5+v_5)$ the exact
coefficient
\begin{equation}\label{eq:cubeBernstein}
\left[\sum_{\ell=0}^{\min(i,\deg c_u)}
 c_{u,\ell}\frac{\binom i\ell}{\binom K\ell}\right]
\prod_{j=1}^5\frac{\binom{s_j}{v_j}}{\binom d{u_j+v_j}}.
\end{equation}
Contributions sharing an index are added. Indeed, after \eqref{eq:stick},
$q^u$ becomes
\[
\prod_{j=1}^5 x_j^{u_j}(1-x_j)^{d-s_j-u_j}.
\]
Multiplication in coordinate $j$ by
$(x_j+(1-x_j))^{s_j}=1$ elevates its degree to $d$ and proves
\eqref{eq:cubeBernstein}. The identity
$\beta^\ell=\sum_{i=\ell}^K\binom i\ell\binom K\ell^{-1}
\binom Ki\beta^i(1-\beta)^{K-i}$ gives the first factor.

Nonnegative Bernstein coefficients imply nonnegativity on the entire cube,
since its basis functions are nonnegative. For an unresolved cube, bisect one
coordinate using midpoint de Casteljau subdivision: from a control line
$b_0,\ldots,b_d$ form
\begin{equation}\label{eq:decasteljau}
b_i^{(0)}=b_i,\qquad
b_i^{(r+1)}=\tfrac12(b_i^{(r)}+b_{i+1}^{(r)}).
\end{equation}
The left child coefficients are $b_0^{(r)}$ for $r=0,\ldots,d$; the right
child coefficient of index $d-r$ is $b_{d-r}^{(r)}$. These are exact
representations on the two half-intervals. Applying this to all parallel
control lines gives the two child tensors. The verifier selects a coordinate
with greatest adjacent-coefficient variation, with ties broken by coordinate
order. All 63 groups terminate in 69 accepted boxes, with depth at most four.

\subsection{The zero-tail group: interval times simplex}

Let $H_0=P_{12}(\beta;q_1,\ldots,q_6,0,\ldots,0)$. Its pure $q_1^{12}$
coefficient is $(1+\beta)^{11}$: setting all later $q$ to zero gives
$m_n=q_1^n$ and $D_{12}(\alpha)=\alpha^{12}q_1^{12}$. Every other composition
monomial contains $q_1q_2$. Thus
\begin{equation}\label{eq:special}
H_0=(1+\beta)^{11}q_1^{12}+q_1q_2R(\beta;q_1,\ldots,q_6),
\end{equation}
where $R$ is homogeneous of degree ten in $q$ and of degree at most ten in
$\beta$. The verifier checks the nonnegative pure coefficient and the common
factor. It remains to certify $R\ge0$ on
$[0,1]\times\{q\ge0:\sum q_i=1\}$.

Write $R=\sum_{|u|=10,\ell}c_{u,\ell}\beta^\ell q^u$. In the product of the
degree-ten interval basis and the degree-ten simplex basis
\[
\binom{10}{i}\beta^i(1-\beta)^{10-i}
\binom{10}{u_1,\ldots,u_6}q^u,
\]
the initial coefficient at $(i,u)$ is
\begin{equation}\label{eq:simplexBernstein}
b_{i,u}=\frac{1}{\binom{10}{u_1,\ldots,u_6}}
\sum_{\ell=0}^{\min(i,10)}
 c_{u,\ell}\frac{\binom i\ell}{\binom{10}\ell}.
\end{equation}
Absent monomials have coefficient zero. There are
$11\binom{15}{5}=33033$ coefficients, in an $11$ by $3003$ array.

A subdivision either bisects the $\beta$ interval, using
\eqref{eq:decasteljau}, or bisects an edge of the simplex. To verify the latter
operation explicitly, choose vertices $a,b$, and fix the other four exponents
$\mu$. The remaining indices have $u_a+u_b=d'=10-|\mu|$. On this line, the
simplex basis is a fixed common factor times the ordinary degree-$d'$
Bernstein basis in $q_a/(q_a+q_b)$, because
\[
\binom{10}{u}q^u
=\frac{10!}{d'!\prod\mu_h!}
\left(\prod q_h^{\mu_h}\right)(q_a+q_b)^{d'}
\binom{d'}{u_a}
\left(\frac{q_a}{q_a+q_b}\right)^{u_a}
\left(\frac{q_b}{q_a+q_b}\right)^{u_b}.
\]
This identity is polynomial after cancellation and extends to $q_a+q_b=0$.
Applying \eqref{eq:decasteljau} to every such control line gives the Bernstein
arrays on the two simplexes obtained by replacing vertex $a$ or $b$ with
their midpoint. Those two simplexes cover the parent simplex; their
intersection is the dividing face. More explicitly, if $\lambda$ denotes
child barycentric coordinates, the child replacing vertex $a$ has
$q_a=\lambda_a/2$, $q_b=\lambda_b+\lambda_a/2$ and
$q_h=\lambda_h$ otherwise; it covers $q_a\le q_b$. The other child has
$q_a=\lambda_a+\lambda_b/2$, $q_b=\lambda_b/2$ and covers $q_b\le q_a$.
These maps prove complete coverage, including the dividing face. This remains
valid after any previous affine subdivision. Each line uses its own degree $d'$, rather than the
ambient degree ten.

The deterministic subdivision choice is included completely in the verifier:
it examines the most negative coefficient and scores the interval line and
the fifteen edge lines through it by the minimum new control coefficient,
then by variation. The choice affects efficiency only. A region is accepted
solely when all exact Bernstein coefficients are nonnegative. The critical
calculation terminates in 76 accepted regions, with depth at most twelve.

\begin{proposition}[Finite polynomial certificate]\label{prop:certificate}
For $q_1,\ldots,q_{12}\ge0$ and $0\le\beta\le1$, one has
$D_{12}(1+\beta;m(q))\ge0$.
\end{proposition}
\begin{proof}
The complete verifier in Appendix~\ref{app:code} reconstructs the polynomial
using \eqref{eq:firstreturn} and \eqref{eq:Bell}. It reads no stored coefficient
table. It converts coefficients by \eqref{eq:cubeBernstein} and
\eqref{eq:simplexBernstein}, and subdivides by \eqref{eq:decasteljau}.
Every arithmetic coefficient and comparison uses integers or Python's exact
rational type \texttt{Fraction}; NumPy supplies object-array containers only.
No floating-point tolerance is used. The verifier raises an exception for
any unresolved or negative region. Its complete results are
\[
\begin{array}{l|rrr}
\text{part}&\text{groups}&\text{accepted regions}&\text{maximum depth}\\\hline
\text{nonzero tails (cubes)}&63&69&4\\
\text{zero tail (interval--simplex)}&1&76&12\\\hline
\text{total}&64&145&12
\end{array}
\]
Every subdivision covers the parent, every terminal region has nonnegative
coefficients, and every group is checked. The pure term in
\eqref{eq:special} and all removed monomial factors are nonnegative. Hence
$P_{12}\ge0$, and multiplication by $1+\beta>0$ proves the proposition.
\end{proof}

\section{Lower derivatives and all real exponents}

\begin{lemma}\label{lem:descending}
Let $g\in C^N(0,\infty)$ be bounded on every ray $[a,\infty)$ and satisfy
$(-1)^Ng^{(N)}\ge0$. Then $(-1)^jg^{(j)}\ge0$ for $1\le j\le N$.
\end{lemma}
\begin{proof}
For $N=1$ the conclusion is already the hypothesis. Otherwise descend through
$j=N-1,\ldots,1$. If $(-1)^{j+1}g^{(j+1)}\ge0$, then
$h=(-1)^jg^{(j)}$ is nonincreasing. Were $h(t_1)<0$ at any point, one would
have $h(t)\le h(t_1)<0$ for $t\ge t_1$. Repeated integration $j$ times
would force $g$ to be unbounded on that ray: the leading term has magnitude
at least a positive constant times $(t-t_1)^j$, whereas the integration
constants have degree at most $j-1$. This is a contradiction. Therefore
$h\ge0$, permitting the next step. The sign of $g$ itself is separate.
\end{proof}

For $1\le\alpha\le2$, Proposition~\ref{prop:certificate},
Lemma~\ref{lem:Stieltjes}, and~\eqref{eq:derivative} prove
$(f^\alpha)^{(12)}\ge0$ at every positive time. Since $f$ is
nonnegative and nonincreasing, $f^\alpha$ is bounded on every such ray.
Lemma~\ref{lem:descending} gives every lower derivative sign, and order zero
is its nonnegativity.

For arbitrary $\alpha\ge1$, write $\alpha=k+\gamma$, where $k$ is a
nonnegative integer and $1\le\gamma\le2$. The function $f^k$ is completely
monotone, by repeated application of Leibniz's rule (with $f^0=1$), while
$f^\gamma\in\CM^{12}$ by the preceding paragraph. The product of two
$\CM^{12}$ functions lies in $\CM^{12}$, since every term of its signed
Leibniz formula is nonnegative. Thus $f^\alpha=f^kf^\gamma\in\CM^{12}$,
proving Theorem~\ref{thm:main}.

\section{An exact obstruction at order thirteen}

Take $f(t)=1+e^{-t}$, $\alpha=6/5$, and $t_0=\log(7/3)>0$.
This $f$ is completely monotone. At $t_0$ its normalized moments satisfy
$m_0=1$ and $m_j=3/10$ for every $j\ge1$. Therefore
\[
\frac{(-1)^{13}(f^{6/5})^{(13)}(t_0)}{f(t_0)^{6/5}}
=\sum_{k=1}^{13}\left\{\!\begin{matrix}13\\k\end{matrix}\!\right\}
\fall{(6/5)}{k}\left(\frac3{10}\right)^k
=-\frac{6524240334632293437}{2980232238769531250}<0.
\]
Here the braces denote Stirling numbers of the second kind. The equality is
an exact rational calculation, reproduced in Appendix~\ref{app:13}.
Consequently, if $N_*$ is the largest integer such that
$f^\alpha\in\CM^{N_*}$ for every $f\in\CM$ and every $\alpha>1$, this note proves
\[
N_*=12.
\]
The finite certificate and analytic reduction together prove this sharp
threshold. The code is an exact computational proof, not a proof-assistant
formalization. Independent external review and a fuller priority review remain
appropriate before publication.

\section{A consequence for discrete Hardy--Rellich--Birman weights}

With the operators and weights of \v Stampach--Waclawek \cite{SW},
Theorem~5.1(ii) extends from $3\le\ell\le7$ to
\[
3\le\ell\le12,\qquad p>2.
\]
Precisely, let $u$ and $\rho$ be the sequence and discrete $p$-polyharmonic
weight defined in their equations (2.16)--(2.17). Then $\rho$ is a discrete
$p$-Hardy--Rellich--Birman weight, is critical, is nonattainable on
$H_0^\ell$, and is optimal near infinity, with these terms understood as in
\cite{SW}. This is a corollary of their theorem and its proof, rather than
an assertion of any stronger optimality property.

Indeed, in Proposition~5.4, the proof labelled ``Ad (ii)'', the restriction
$\ell\le7$ comes from the inclusion that the $(p-1)$st power of a completely
monotone function belongs to $\CM^\ell$, quoted in their Lemma~5.3(2).
Here $p-1>1$, so Theorem~\ref{thm:main} supplies that inclusion for every
$\ell\le12$. The remaining arguments in their Proposition~5.4 and
Theorem~5.1(ii) apply with this substitution. No claim is made here for the
other parameter cases of their theorem.

\begin{thebibliography}{9}
\bibitem{vH} H. van Haeringen, \emph{Inequalities for real powers of completely
monotonic functions}, Journal of Mathematical Analysis and Applications
\textbf{210} (1997), 102--113.
\url{https://doi.org/10.1006/jmaa.1997.5376}.
\bibitem{SW} F. \v Stampach and J. Waclawek,
\emph{Optimal discrete $p$-Hardy--Rellich--Birman inequalities},
arXiv:2605.25238 (2026), Lemma~5.3.
\url{https://arxiv.org/html/2605.25238v1}.
\bibitem{PSZ} M. P\'etr\'eolle, A. D. Sokal, and B.-X. Zhu,
\emph{Lattice paths and branched continued fractions: An infinite sequence of
generalizations of the Stieltjes--Rogers and Thron--Rogers polynomials, with
coefficientwise Hankel-total positivity}, arXiv:1807.03271v2 (2020),
Theorem~1.1.
\url{https://arxiv.org/abs/1807.03271}.
\bibitem{Flajolet} P. Flajolet, \emph{Combinatorial aspects of continued
fractions}, Discrete Mathematics \textbf{32} (1980), 125--161.
\url{https://algo.inria.fr/flajolet/Publications/Flajolet80b.pdf}.
\bibitem{Bernstein} R. L. Schilling, R. Song, and Z. Vondra\v cek,
\emph{Bernstein Functions}, second edition, De Gruyter, 2012,
Theorem~1.4 (Bernstein's representation theorem).
\end{thebibliography}

\appendix
\section{Complete exact order-twelve verifier}\label{app:code}
This standalone Python 3 program requires NumPy, used solely for object arrays.
All proof arithmetic is integer or \texttt{Fraction} arithmetic. It reconstructs
all coefficients, uses no numerical acceptance tolerance, and rejects an
unresolved region. Copy the listing to \texttt{verify\_cm12.py} and run
\texttt{python verify\_cm12.py}; the output ends with 64 groups, 145 leaves,
and maximum depth twelve. Timing measurements do not enter the proof.
\begin{lstlisting}
#!/usr/bin/env python3
"""Exact certificate: positive real powers preserve CM derivatives through 12.

Every coefficient and subdivision operation is a Fraction or integer. NumPy is
used only as an object-array container; no floating-point operations or acceptance
tolerances occur. The accompanying proof gives the analytic moment reduction,
Bernstein geometry, lower-derivative implication, and extension in alpha.
"""
import math
from itertools import combinations
from fractions import Fraction
from fractions import Fraction as F
from itertools import product
from collections import defaultdict
from functools import lru_cache
from math import comb
import numpy as np
import sys,time

NMAX=12
zero=(0,)*NMAX
def mul(a,b):
 c=defaultdict(int)
 for p,v in a.items():
  for q,w in b.items():c[tuple(x+y for x,y in zip(p,q))]+=v*w
 return dict(c)
@lru_cache(None)
def moment(h,n):
 if n==0:return {zero:1}
 ans=defaultdict(int)
 for j in range(n):
  for p,v in mul(moment(h,j),moment(h+1,n-1-j)).items():
   q=list(p);q[h]+=1;ans[tuple(q)]+=v
 return dict(ans)
@lru_cache(None)
def bell(n,k):
 if n==k==0:return {zero:1}
 if n==0 or k==0:return {}
 ans=defaultdict(int)
 for j in range(1,n-k+2):
  for p,v in mul(moment(0,j),bell(n-j,k-1)).items():ans[p]+=comb(n-1,j-1)*v
 return dict(ans)
def pmul(a,b):
 c=[0]*(len(a)+len(b)-1)
 for i,x in enumerate(a):
  for j,y in enumerate(b):c[i+j]+=x*y
 return c
fall=[[1]]
for k in range(1,NMAX+1):fall.append(pmul(fall[-1],[2-k,1]))
def derivative(n):
 ans=defaultdict(lambda:[0]*(n+1))
 for k in range(1,n+1):
  for p,v in bell(n,k).items():
   for j,w in enumerate(fall[k]):ans[p][j]+=v*w
 return ans
def div_alpha(c):
 q=[c[0]]
 for k in range(1,len(c)-1):q.append(c[k]-q[-1])
 if c[-1]!=q[-1]:raise ArithmeticError('alpha division')
 while len(q)>1 and q[-1]==0:q.pop()
 return q
def split(a,axis):
 x=np.moveaxis(a,axis,0);n=x.shape[0]-1
 l=np.empty_like(x);r=np.empty_like(x)
 l[0]=x[0];r[n]=x[n]
 for k in range(1,n+1):
  x=(x[:-1]+x[1:])/F(2)
  l[k]=x[0];r[n-k]=x[-1]
 return np.moveaxis(l,0,axis),np.moveaxis(r,0,axis)
def check(a,depth=0,maxdepth=12):
 if min(a.flat)>=0:return True,1,depth
 if max(a.flat)<0:return False,1,depth
 if depth>=maxdepth:return None,1,depth
 varies=[max(abs(v) for v in np.diff(a,axis=j).flat) if a.shape[j]>1 else F(0) for j in range(a.ndim)]
 axis=max(range(a.ndim),key=lambda j:varies[j])
 c,d=split(a,axis)
 ok1,n1,d1=check(c,depth+1,maxdepth)
 if ok1 is not True:return ok1,n1,d1
 ok2,n2,d2=check(d,depth+1,maxdepth)
 return ok2,n1+n2,max(d1,d2)
def certify_nonzero_tails(n=12,r=6):
 groups=defaultdict(dict)
 coeff=derivative(n)
 if len(coeff)!=2**(n-1):raise ArithmeticError('monomial coverage')
 for p,c in coeff.items():groups[p[r:]][p[:r]]=div_alpha(c)
 if len(groups)!=64:raise ArithmeticError('group count')
 leaves=0;depth=0;fail=[];checked=0
 for tail,poly in groups.items():
  if not any(tail):continue
  checked+=1
  mins=tuple(min(e[j] for e in poly) for j in range(r))
  D=n-sum(tail)-sum(mins);K=max(map(len,poly.values()))-1
  arr=np.full((K+1,)+(D+1,)*(r-1),F(0),dtype=object)
  for e,c in poly.items():
   u=[e[j]-mins[j] for j in range(r)]
   prefixes=[sum(u[:j]) for j in range(r-1)]
   for vs in product(*(range(p+1) for p in prefixes)):
    inds=tuple(u[j]+vs[j] for j in range(r-1))
    weight=F(1)
    for j in range(r-1):weight*=F(comb(prefixes[j],vs[j]),comb(D,inds[j]))
    for i in range(K+1):
     b=sum((F(c[j]*comb(i,j),comb(K,j)) for j in range(min(i,len(c)-1)+1)),F(0))
     arr[(i,)+inds]+=b*weight
  ok,nn,dd=check(arr);leaves+=nn;depth=max(depth,dd)
  if ok is not True:fail.append((tail,ok,nn,dd))
 if checked!=63:raise ArithmeticError('nonzero-tail coverage')
 print('EXACT',n,'prefix',r,'nonzero-tail groups',checked,'leaves',leaves,'maxdepth',depth,'failed',len(fail),flush=True)
 if fail:print(fail[:10],flush=True)
 if fail:raise ArithmeticError('nonzero-tail certificate failed')
 return leaves,depth

@lru_cache(None)
def weak(n,r):
 if r==1:return ((n,),)
 return tuple((i,)+p for i in range(n+1) for p in weak(n-i,r-1))

def certify_zero_tail(n=12,r=6,maxdepth=25):
 poly={p[:r]:div_alpha(c) for p,c in derivative(n).items() if not any(p[r:])}
 pure=(n,)+(0,)*(r-1)
 if min(poly[pure])<0:raise ArithmeticError('pure negative')
 del poly[pure]
 mins=tuple(min(p[j] for p in poly) for j in range(r))
 if mins!=(1,1,0,0,0,0):raise ArithmeticError('special common monomial')
 D=n-sum(mins);K=max(map(len,poly.values()))-1
 exps=weak(D,r);loc={p:i for i,p in enumerate(exps)}
 arr=np.full((K+1,len(exps)),Fraction(0),dtype=object)
 for p,c in poly.items():
  u=tuple(p[j]-mins[j] for j in range(r));mult=math.factorial(D)
  mult=math.factorial(D)//math.prod(math.factorial(uj) for uj in u)
  for i in range(K+1):arr[i,loc[u]]=sum((Fraction(c[j]*math.comb(i,j),math.comb(K,j)) for j in range(min(i,len(c)-1)+1)),Fraction(0))/mult
 if not all(isinstance(v,Fraction) for v in arr.flat):raise ArithmeticError('exact coefficient types')
 # Exact arithmetic: no normalization or tolerance.
 lines={};lineof={}
 for i,j in combinations(range(r),2):
  gs=defaultdict(list)
  for p in exps:gs[tuple(p[h] for h in range(r) if h not in(i,j))].append(p)
  lines[(i,j)]=[np.array([loc[p] for p in sorted(g,key=lambda q:q[i])]) for g in gs.values()]
  lineof[(i,j)]={int(v):line for line in lines[(i,j)] for v in line}
 print('shape',arr.shape,'D,K',D,K,'mins',mins,flush=True)
 nodes=0
 def split_beta(a):
  x=a.copy();left=np.zeros_like(a);right=np.zeros_like(a);left[0]=x[0];right[K]=x[-1]
  for kk in range(1,K+1):
   x=(x[:-1]+x[1:])/Fraction(2);left[kk]=x[0];right[K-kk]=x[-1]
  return left,right
 def split_edge(a,edge):
  left=np.zeros_like(a);right=np.zeros_like(a)
  for inds in lines[edge]:
   x=a[:,inds];d=len(inds)-1
   left[:,inds[0]]=x[:,0];right[:,inds[d]]=x[:,-1]
   for kk in range(1,d+1):
    x=(x[:,:-1]+x[:,1:])/Fraction(2)
    left[:,inds[kk]]=x[:,0];right[:,inds[d-kk]]=x[:,-1]
  return left,right
 def check(a,depth=0,path=()):
  nonlocal nodes
  nodes+=1
  if min(a.flat)>=0:return True,1,depth
  if max(a.flat)<0:return False,1,depth
  if depth>=maxdepth:
   print('FAILPATH',path,'range',a.min(),a.max(),flush=True)
   return None,1,depth
  ib,jb=np.unravel_index(a.argmin(),a.shape)
  def score(values):
   values=values.copy();mn=min(values[0],values[-1]);variation=max(abs(np.diff(values)),default=0)
   for kk in range(1,len(values)):
    values=(values[:-1]+values[1:])/Fraction(2);mn=min(mn,values[0],values[-1])
   return mn,variation
  scores={None:score(a[:,jb])}
  for edge in lines:scores[edge]=score(a[ib,lineof[edge][jb]])
  edge=max(scores,key=scores.get)
  le,ri=split_beta(a) if edge is None else split_edge(a,edge)
  o1,n1,d1=check(le,depth+1,path+((edge,0),))
  if o1 is not True:return o1,n1,d1
  o2,n2,d2=check(ri,depth+1,path+((edge,1),))
  return o2,n1+n2,max(d1,d2)
 start=time.time();initial=[arr]
 for lev in range(0):initial=[v for a in initial for v in split_beta(a)]
 results=[]
 for idx,a in enumerate(initial):
  ans=check(a);results.append(ans)
  print('SLICE',idx,'RESULT',ans,'nodes',nodes,'seconds',time.time()-start,flush=True)
  if ans[0] is not True:break
 if any(ans[0] is not True for ans in results):raise ArithmeticError('zero-tail certificate failed')
 print('ALLRESULTS',results,flush=True)
 return sum(ans[1] for ans in results),max(ans[2] for ans in results)

if __name__=='__main__':
 start=time.time()
 ordinary=certify_nonzero_tails()
 special=certify_zero_tail()
 print('PASS: all 64 groups at order 12 have exact nonnegative Bernstein certificates.')
 print('Total leaves:',ordinary[0]+special[0],'; maximum depth:',max(ordinary[1],special[1]))
 print('Elapsed seconds:',time.time()-start)
\end{lstlisting}

\section{Exact order-thirteen witness}\label{app:13}
The following standard-library-only script evaluates the rational obstruction.
\begin{lstlisting}
#!/usr/bin/env python3
"""Exact thirteenth-derivative obstruction for powers of CM functions.

f(t) = 1 + exp(-t), alpha = 6/5, t0 = log(7/3).
At t0 the normalized derivative moments of f are m_j = 3/10 (j>=1).
The normalized signed derivative of f**alpha is therefore the finite
Stirling polynomial computed below. No numerical transcendental evaluation
is needed to decide its sign.
"""
from fractions import Fraction as Q
import json


def certificate():
    n = 13
    alpha, p = Q(6, 5), Q(3, 10)
    stirling = [[0] * (n + 1) for _ in range(n + 1)]
    stirling[0][0] = 1
    for j in range(1, n + 1):
        for k in range(1, j + 1):
            stirling[j][k] = stirling[j-1][k-1] + k*stirling[j-1][k]
    falling, total = Q(1), Q(0)
    for k in range(1, n + 1):
        falling *= alpha-k+1
        total += stirling[n][k] * falling * p**k
    expected = Q(-6524240334632293437, 2980232238769531250)
    if total != expected or total >= 0:
        raise RuntimeError("Exact derivative obstruction failed")
    return {"f": "1 + exp(-t)", "alpha": str(alpha),
            "time": "log(7/3)", "derivative_order": n,
            "normalized_signed_derivative": str(total),
            "strictly_negative": True}


if __name__ == "__main__":
    print(json.dumps(certificate(), indent=2))
\end{lstlisting}
\end{document}
